% Combined midterm report. Compile from submission/report/:
%   latexmk -pdf main.tex
% Course limit: 5 pages including figures. References may continue past that.

\documentclass[10pt,twocolumn,letterpaper]{article}
\usepackage[T1]{fontenc}

\usepackage[pagenumbers]{cvpr}
%% This file contains a number of tweaks that are typically applied to the main document.
%% They are not enabled by default, but can be enabled by uncommenting the relevant lines.

%%
%% Packages used by this report that cvpr.sty does not already load
%%
\usepackage{multirow}

%%
%% Inline annotations; for predefined colors, refer to "dvipsnames" in the xcolor package:
%% https://tinyurl.com/overleaf-colors
%%
\newcommand{\red}[1]{{\color{red}#1}}
\newcommand{\todo}[1]{{\color{red}#1}}
\newcommand{\TODO}[1]{\textbf{\color{red}[TODO: #1]}}
%%
%% disable for camera ready / submission by uncommenting these lines  
%%
% \renewcommand{\TODO}[1]{}
% \renewcommand{\todo}[1]{#1}

%%
%% work harder in optimizing text layout. Typically shrinks text by 1/6 of page, enable
%% it at the very end of the writing process, when you are just above the page limit
%%
\usepackage{microtype}

%%
%% fine-tune paragraph spacing
%%
% \renewcommand{\paragraph}[1]{\vspace{.5em}\noindent\textbf{#1.}}

%%
%% globally adjusts space between figure and caption
%%
% \setlength{\abovecaptionskip}{.5em}


%%
%% Allows "the use of \paper to refer to the project name"
%% with automatic management of space at the end of the word
%%
% \usepackage{xspace}
% \newcommand{\paper}{ProjectName\xspace}

%%
%% Commonly used math definitions
%%
% \DeclareMathOperator*{\argmin}{arg\,min}
% \DeclareMathOperator*{\argmax}{arg\,max}

%%
%% Tigthen underline
%%
% \usepackage{soul}
% \setuldepth{foobar}

\definecolor{cvprblue}{rgb}{0.21,0.49,0.74}
\usepackage[pagebackref,breaklinks,colorlinks,allcolors=cvprblue]{hyperref}

\def\paperID{*****}
\def\confName{CVPR}
\def\confYear{2026}

\title{Semi-Supervised Labels and Outcome Prediction\\
for Clinical Trials\\
\vspace{0.35em}
\large CS5785 / ORIE5750 / ECE5414 --- Midterm Project}

\author{Anirudha Hadagali\\
Cornell University
\and
Michael Mazzola\\
Cornell University
\and
Tarun Mandalapu\\
Cornell University
}

\begin{document}
\maketitle

\begin{abstract}
We address two clinical trial prediction tasks. In Part~1, only 157 of 1,171
training trials have observed stopping-reason labels. Using 150-dimensional
embeddings, we combine an anchored self-training classifier with a second
pipeline that adds graph, cluster, and neighbor information. Our fixed 80/20
blend achieved a Kaggle sample-wise F1 of 0.89935, compared with 0.89636 for
anchored alone.

In Part~2, we predict whether a trial meets its primary endpoint. To account
for the later test period, we give recent trials more weight and incorporate
neighbors that favor recent matches. Blending in 30\% of an earlier
text-and-tree ensemble achieved a Kaggle log loss of 0.58465. A later blend
improved our historical evaluation scores but scored 0.58492 on Kaggle, so we
retained the earlier model.

\end{abstract}

\noindent\textbf{Team responsibilities.}
All three authors worked on both parts and on this report.
\textbf{Michael Mazzola} led the Part~1 model.
\textbf{Tarun Mandalapu} built the first Part~2 classifier, and
\textbf{Anirudha Hadagali} led the later Part~2 tuning.
Tarun Mandalapu also drafted the project notes used to prepare this report.

\section{Methods}
\label{sec:methods}

\subsection{Part 1: stopping reasons}
\label{sec:p1}

\subsubsection{Data and evaluation}

The dataset contains 1,561 trials that stopped early~\cite{fu2022hint,clinicaltrials}.
The training set has 1,171 trials: 157 with observed labels and 1,014 with every
reason marked $-100$, indicating an unknown label. The validation set has 78
labeled trials, and the test set has 312 later trials from 2010 to 2016.
A trial can have more than one of five stopping reasons: Safety, Funding,
Efficacy, Recruitment, and Other. Among the 235 labeled trials, Recruitment
appears on 64\% and Funding on 5\%. Most labeled trials are from earlier years,
so these rates may differ in the test set.

We initially selected models and features using sample-wise F1 on the supplied
validation set, alongside per-reason F1 and its macro average. As we tested
more configurations, we became concerned that repeatedly tuning against the
same 78 trials could overfit our model choices to that particular split.
A single trial can change sample-wise F1 by about 1.3 percentage points.
This is selection overfitting, even when validation labels are excluded from
training.

We therefore moved to repeated grouped cross-validation as a more stable guide
for tuning, averaging over eight repeats of five folds on both the 157 training
labels and all 235 labels, called Train157 and Pooled235 in the notebooks and results. This lets
us judge whether a change helps across
several held-out groups rather than only the original validation set. Trials
with identical embeddings stay together, and we remove masked duplicates of
held-out trials. We also use three chronological windows totaling 98
predictions to check performance on later trials. These reused development
splits reduce dependence on a single holdout, but do not eliminate selection
bias or provide a new independent test set.

\subsubsection{Models and training}

We use the 150 supplied stopping-reason embedding coordinates as features.
Trial IDs align rows, and start years define our chronological evaluation
splits. In an earlier self-training model, adding start year, eligibility text,
disease, or drug features~\cite{rogers2010extended,rdkit,milacic2024reactome}
did not improve validation sample-wise F1. Embeddings alone scored .893,
compared with .891 with year, .876 with eligibility or disease text, and .889
with fingerprints. These experiments guided our feature choices, although we
did not repeat the same feature comparisons with the final blend.

Our \emph{anchored} component standardizes embeddings using the labeled and
masked training rows. A 30-component diagonal Gaussian mixture uses 30
principal components~\cite{dempster1977em,pedregosa2011scikit}, with $k$-means
initialization, three restarts, and seed zero. EM stops when the average
lower-bound gain falls below $10^{-3}$, with a maximum of 100 iterations; we
check convergence before using the fit. The logistic classifier allows 3,000
iterations with tolerance $10^{-4}$. For each component, we estimate
reason frequencies from the current training labels, adding one pseudo-count based on the
overall rates. The pseudo-label teacher mixes 5\% of these profiles with 95\%
of the classifier score. Our earlier $k$-means and DBSCAN models performed
worse on Kaggle. DBSCAN-plus-logistic
reached .910 validation F1 under an earlier protocol, but its original and
235-label submissions scored .87500 and .86452 respectively.

The anchored classifier combines multiclass logistic regression ($C=.1$),
cosine 10-nearest-neighbor votes, and shrinkage linear discriminant analysis
(LDA), with fixed weights .5/.3/.2. All use the 150 standardized coordinates.
A multi-label trial is duplicated across its positive classes for the
multiclass fits. Logistic weights divide its mass across positives and use
class weights fixed from observed labels. Neighbors vote for every positive label,
weighted by nonnegative cosine similarity squared. LDA uses equal class priors.

Two pseudo-label refreshes give three classifier fits in total~\cite{yarowsky1995unsupervised,triguero2015selflabeled}.
We accept a masked trial when its highest teacher score is at least .8. At each
refresh, we rebuild the training set from observed labels and the currently
accepted pseudo-labels. These pseudo-labeled rows receive logistic weight .25;
the neighbor and LDA components use them without
that downweighting. The final 235-label anchored fit accepts 757 masked rows in
the last refresh. We keep observed labels fixed throughout self-training.

Our graph-based label-combination model, called \emph{Cand2} in the notebooks,
models eight fixed label combinations. It adds soft eight-cluster $k$-means
membership and 20-neighbor label-vote features
to standardized embeddings. Its Correct-and-Smooth
labeler~\cite{huang2020correctsmooth} uses a symmetric 10-neighbor graph on
scaled, L2-normalized 64-dimensional PCA coordinates. Correction and smoothing
coefficients are .8 and .419962. Generated labels require confidence .5 and
receive weight .0630725 in balanced, weighted shrinkage LDA. To prevent leakage
during evaluation, we refit preprocessing, graph labeling, and calibration within inner grouped
partitions and remove duplicates of held-out trials. We fit temperature
calibration using the inner held-out predictions. We average five seeds and
convert combination probabilities into five reason marginals.

The selected rule combines anchored raw scores and Cand2 marginals:
\[
s = .8s_{\mathrm{anchored}} + .2p_{\mathrm{Cand2}}.
\]
We output every reason with $s_j\geq .3$ and always include the highest-scoring
reason. We treat these scores as decision scores rather than calibrated
probabilities and keep the blend weights fixed. For standalone
Cand2, we instead decode label combinations by maximizing expected F1.

Beyond these components, we compared known-label logistic regression,
nearest-neighbor classifiers, Gaussian naive Bayes, RBF support vector
machines, and LDA ensembles. We tested $k$-means, GMM, DBSCAN, and graph-based
label generation, then explored kernel and symbolic-regression
(PySR~\cite{cranmer2023pysr}) blend partners,
learned combiners, and alternative treatment of secondary Other labels. These
comparisons cover different ways to generate labels, represent trials, and
combine predictions; the notebooks retain the full configurations.

We fit both components using training inputs only and use test embeddings only
for prediction. For our final submission, we refit anchored on all 235 observed
labels plus the 1,014 masked rows and combine its scores with the saved
five-seed Cand2 predictions. The .891 validation score comes from a separate
fit using only the 157 training labels; it does not independently evaluate the
final model trained on all 235 labels. Our final-model notebook documents both
components and the complete blend.
For local blend comparisons, we combine the components' saved held-out scores
on the same trials, without retraining them for each set of blend weights.

\subsection{Part 2: trial outcomes}
\label{sec:p2}

\subsubsection{Data and evaluation}

Each row represents one trial in one phase. Combined-phase trials appear twice,
so we keep their rows together in every split. The training set has 8,037 rows
through 2013, validation has 1,013 rows from the same years, and the test set
has 1,500 later rows from 2013 to 2022, with 500 per phase. We evaluate models
using the average of the three phase log losses, matching the Kaggle metric.
We also examine F1 at a 0.5 cutoff, ROC-AUC, and precision-recall AUC, computed as average precision with success as the positive
class. We join Part~1 stopping reasons after prediction to analyze errors;
neither reason labels nor reason embeddings enter the outcome models.

\subsubsection{Models and training}

Our baseline uses logistic regression on eligibility text: TF-IDF features from
words and bigrams, at most 35,000 terms, minimum document frequency 3,
sublinear term frequency,
three phase indicators, and $C = 2$. Separate phase-only fits overfit Phase~I.
Drug fingerprints did not reliably help, and a matched drug is less common in
test (39\%) than in training (46\%). Naive Bayes ranked trials reasonably well
but produced overconfident probabilities, resulting in poor log loss.

We then built on a stack of eight classifiers~\cite{wolpert1992stacked}:
logistic models on eligibility text, diseases, drug names, codes, drug
annotations, and metadata, plus naive Bayes and a 50-neighbor model on the
eligibility text. A second logistic regression learns to combine them from
five-fold out-of-fold predictions, keeping each trial within a single fold.
Recent trials get more weight, with a three-year exponential time constant, so
a trial loses about half its weight every two years. That stack, refit on all
9,050 labeled rows, scored 0.58688. The same stack without recency weights
scored 0.58810. Our first ensembles had used only the 8,037 training rows, so
part of the earlier gap was the extra data.

The neighbor model within the stack ignored time and pulled estimates toward a
50\% success rate. To address this, we created a recency-and-neighbor blend
(submission file~10 in the notebooks) by averaging the recency
stack with a second stack that uses a shorter, two-year time constant. The
second stack also includes eight neighborhood estimates that vary the number
of neighbors and similarity weighting, with some incorporating disease, drug,
and ICD text. We give the two stacks equal weight. This improved every
chronological window and achieved 0.58609 on Kaggle, although log loss on the
random validation split became slightly worse.

File~10 performed less well on splits that mixed years, while our earlier
text-and-tree ensemble was weaker on later trials. Its trees use
CatBoost~\cite{prokhorenkova2017catboost}, with 1,000 boosting iterations,
learning rate .035, and L2 leaf regularization 6.
The earlier ensemble assigns
half its weight to three logistic models, a quarter to a depth-6 tree, and a
quarter to a depth-8 tree. Combining the models helped: a 20\% contribution
from the earlier ensemble achieved 0.58474 on Kaggle. Increasing its share to
30\% reduced log loss by a further 0.00009. We selected this blend, file~13,
with submitted probability
\[
0.70\,P_{10} + 0.30\,P_{\text{early}},
\quad
P_{10} = 0.50\,P_{\text{recency}} + 0.50\,P_{\text{fast neighbors}}.
\]
We clip the final probabilities to $[.001,.999]$.
After fixing the model configuration, we refit the supervised components on all
9,050 labeled rows. The original stack allows vocabulary construction to use
unlabeled test text, but no test labels are used in training. A later blend (file~15)
with an additional tree and a phase trend improved our local scores but
achieved a worse Kaggle log loss of 0.58492, so we retained file~13.

\subsubsection{Validation and error analysis}

For later outcome candidates, we compare the supplied validation split,
two repeats of three-fold trial-grouped cross-validation, and three forward
windows trained through 2007, 2009, or 2011. The grouped checks compare models
across held-out trials, while the forward windows assess the shift to later
test years. We refit preprocessing within each training
split, including inner stacking folds, and track the worst temporal window as
well as the mean. All official validation metrics below hold out the 1,013
validation outcomes from fitting, even though we later refit the submission
model on all 9,050 labeled rows.

For failed validation trials, we match trial IDs to Part~1. We use observed
training or validation reason labels where available; otherwise we use
predictions from our frozen historical DBSCAN/logistic model, with a 0.5 cutoff
and highest-score fallback. We keep this reason assignment fixed across outcome
models. Unmatched trials form the ``no recorded stopping reason'' group.
Multi-label trials enter each applicable category, and a combined-phase trial
is counted once within each phase. The false-success rate is the fraction of
failed trials assigned success probability at least 0.5.

\section{Results}
\label{sec:results}

\subsection{Part 1: stopping reasons}

Table~\ref{tab:p1} compares our Part~1 models on the same grouped splits.
We tested small contributions from LDA or Cand2 while keeping anchored as the
main component. Cand2 was weaker on its own, but its graph and label-combination
approach could complement anchored's predictions. The 80/20 blend matched
anchored's original validation F1 and remained close on both CV pools, with
small gains in their means and the pooled chronological score.

We were concerned that anchored might not generalize as well to later trials.
Cand2's original search considered CV and chronological performance through
the combined score $J$. It outperformed anchored on $J$'s chronological checks,
although this advantage did not hold on our stricter chronology benchmark.
That mixed evidence motivated trying a small Cand2 contribution while keeping
most of anchored's predictions. The blend maintained similar local performance
and subsequently scored .89935 on Kaggle, compared with .89850 for
anchored/LDA, so we retained it as our final model.

On the 78-row validation set, the selected blend predicts the same labels as
anchored. Its per-reason F1 is .941 for Safety (9 positives), .667 for Funding
(4), .759 for Efficacy (14), .990 for Recruitment (50), and .471 for Other (10).
Overall, sample-wise F1 is .891 and macro F1 is .765. These validation results
reflect our development split; the Kaggle column gives the leaderboard scores
recorded for our submissions. The repeated splits overlap, so these results
do not establish statistical significance.

\begin{table}[t]
\centering\small
\setlength{\tabcolsep}{2pt}
\begin{tabular*}{\columnwidth}{@{\extracolsep{\fill}}lccccc@{}}
\toprule
Model & Train & Pooled & Val. & Time & Kaggle \\
\midrule
80/20 Cand2 & \textbf{.9352} & .9216 & .891 & \textbf{.888} & \textbf{.89935} \\
75/25 LDA & .9321 & \textbf{.9234} & .891 & .884 & .89850 \\
Anchored & .9321 & .9214 & .891 & .884 & .89636 \\
LDA ensemble & .9318 & .9216 & \textbf{.906} & .878 & .89401 \\
20\% PySR & .9275 & .9198 & .885 & .878 & .89423 \\
20\% kernel & .9334 & .9209 & .887 & .871 & .89316 \\
\bottomrule
\end{tabular*}
\caption{Part~1 sample-wise F1, higher is better. Train and pooled average
eight grouped repeats on 157 and 235 labels; validation uses 78 trials; time
pools three windows. Cand2, LDA, and kernel rows blend with anchored. PySR's
parent is the earlier 75/25 anchored/LDA blend.}
\label{tab:p1}
\end{table}

\subsection{Part 2: outcome prediction}

Table~\ref{tab:p2} summarizes our Part~2 results; lower log loss is better.
We selected file~13. File~15 improved both local scores but performed worse
on Kaggle, showing that these local gains did not carry over.

\begin{table}[t]
\centering\small
\setlength{\tabcolsep}{3pt}
\begin{tabular*}{\columnwidth}{@{\extracolsep{\fill}}lccc@{}}
\toprule
Model & Time & Grouped & Kaggle \\
\midrule
Phase-specific ensemble & --- & --- & 0.59562 \\
Text-heavy ensemble & --- & --- & 0.58993 \\
Stack, no recency & 0.5421 & 0.5797 & 0.58810 \\
Recency stack & 0.5374 & 0.5874 & 0.58688 \\
File 10, recent neighbors & 0.5359 & 0.5869 & 0.58609 \\
File 13, 30\% early ensemble & 0.5334 & 0.5790 & \textbf{0.58465} \\
File 15, tree and trend & \textbf{0.5327} & \textbf{0.5783} & 0.58492 \\
\bottomrule
\end{tabular*}
\caption{Part~2 log loss. Time averages the forward windows. Grouped is
trial-grouped cross-validation. The first two submissions used an earlier
local protocol, so those cells are blank.}
\label{tab:p2}
\end{table}

\subsubsection{Validation by phase}

Table~\ref{tab:phase} compares the same held-out validation rows. Relative to
the baseline, our final model improves ROC-AUC in all three phases and F1 in
Phases~I and~III, while Phase~II F1 falls from .657 to .647. Phase~III log loss
improves from .529 to .513, but Phase~I rises from .556 to .559. Thus the
modification has mixed effects even though the phase-averaged log loss improves.
The final model's phase-averaged validation log loss is .55328, compared with
.58465 on Kaggle, a gap of .03137 consistent with the later test cohort being
harder to predict. File~10 also illustrates the selection trade-off: it is
weaker than the recency stack on random validation but stronger on the forward
windows and Kaggle.

\begin{table}[t]
\centering\small
\setlength{\tabcolsep}{2pt}
\begin{tabular}{llcccc}
\toprule
Phase & Metric & Base & Recency & File 10 & Final \\
\midrule
\multirow{4}{*}{I}
 & F1 & 0.765 & \textbf{0.793} & 0.762 & 0.788 \\
 & ROC-AUC & 0.783 & \textbf{0.794} & 0.787 & 0.791 \\
 & PR-AUC & \textbf{0.835} & 0.830 & 0.824 & 0.832 \\
 & log loss & 0.556 & 0.556 & 0.560 & 0.559 \\
\midrule
\multirow{4}{*}{II}
 & F1 & \textbf{0.657} & 0.647 & 0.638 & 0.647 \\
 & ROC-AUC & 0.740 & 0.753 & 0.754 & \textbf{0.757} \\
 & PR-AUC & 0.717 & 0.732 & 0.730 & \textbf{0.735} \\
 & log loss & 0.600 & 0.593 & 0.593 & \textbf{0.587} \\
\midrule
\multirow{4}{*}{III}
 & F1 & 0.833 & 0.825 & 0.829 & \textbf{0.839} \\
 & ROC-AUC & 0.774 & 0.790 & 0.786 & \textbf{0.795} \\
 & PR-AUC & 0.862 & 0.882 & 0.877 & \textbf{0.883} \\
 & log loss & 0.529 & 0.516 & 0.521 & \textbf{0.513} \\
\bottomrule
\end{tabular}
\caption{Part~2 held-out validation metrics. Final is file~13; Base is the
eligibility logistic model. F1 uses a 0.5 cutoff; PR-AUC is average precision.
Phase sizes are 149, 501, and 363. Best value per row is bold; ties at the
displayed precision are unbolded.}
\label{tab:phase}
\end{table}

\subsubsection{Errors by stopping reason}

Table~\ref{tab:err} compares the baseline and final model using the same reason
groups. False-success rates fall for Phase~I Recruitment (.36 to .09) and
Phase~II Recruitment (.20 to .14), and for trials without a recorded reason in
Phases~I and~II. Phase~III remains difficult: Recruitment errors increase from
.52 to .57, while the no-reason rate stays at .56. Efficacy errors fall from
8/9 to 7/9, but all nine reason tags are predictions. Small groups also move
in opposite directions: Phase~I Safety improves from 5/9 to 3/9, while Efficacy
worsens from 1/2 to 2/2. These small counts make category-specific gains
uncertain.

\begin{table}[t]
\centering\small
\setlength{\tabcolsep}{2.5pt}
\begin{tabular*}{\columnwidth}{@{\extracolsep{\fill}}llcccc@{}}
\toprule
Phase & Reason & $n$ & Obs./Pred. & Base & Final \\
\midrule
I & Safety & 9 & 1/8 & .56 & .33 \\
I & Funding & 4 & 1/3 & .50 & .25 \\
I & Efficacy & 2 & 0/2 & .50 & 1.00 \\
I & Recruitment & 22 & 2/20 & .36 & .09 \\
I & Other & 5 & 0/5 & .80 & .40 \\
I & No reason & 28 & --- & .43 & .32 \\
II & Safety & 15 & 1/14 & .40 & .33 \\
II & Funding & 4 & 1/3 & .00 & .00 \\
II & Efficacy & 10 & 0/10 & .30 & .30 \\
II & Recruitment & 87 & 12/75 & .20 & .14 \\
II & Other & 18 & 3/15 & .33 & .22 \\
II & No reason & 144 & --- & .35 & .28 \\
III & Safety & 6 & 1/5 & .50 & .50 \\
III & Funding & 0 & 0/0 & --- & --- \\
III & Efficacy & 9 & 0/9 & .89 & .78 \\
III & Recruitment & 21 & 5/16 & .52 & .57 \\
III & Other & 8 & 2/6 & .63 & .63 \\
III & No reason & 82 & --- & .56 & .56 \\
\bottomrule
\end{tabular*}
\caption{False-success rates among failed validation trials. Obs./Pred.
counts observed and predicted reason tags within each group. ``No reason''
means absent from Part~1. Multi-label trials enter each applicable reason.}
\label{tab:err}
\end{table}

\section{Discussion}
\label{sec:why}

\subsection{Part 1: what changed our approach}

We began with logistic regression trained only on the observed labels. This
gave us a reference for deciding whether semi-supervised learning actually
helped, rather than assuming that adding generated labels would improve the
model. The first $k$-means experiments showed the risk: their generated labels
mostly reinforced Recruitment and hurt validation performance, especially for
rarer reasons. We learned to pay attention to which labels we were adding,
not just how many trials we could add. This motivated more conservative label
generation and the combination of supervised scores with clustering.

Our evaluation procedure changed as we experimented. Initially, we relied
heavily on the supplied validation set, but repeatedly choosing settings on
just 78 trials could favor quirks of that split. We moved to repeated CV on
both the 157 training labels and the pooled 235 training-plus-validation
labels. The first kept the original training setting as a reference; the
second let each fold fit on more observed labels while still holding out the
trials being evaluated. Looking for gains across both pools gave us a more
stable basis for tuning than repeatedly chasing the best validation score.
These were still development comparisons, not independent confirmation.

We also tested whether the available clinical and drug information could add
to the embeddings. We added eligibility text, disease and drug information,
pathways, molecular features, and year to the strongest models at that stage,
but the feature screens did not improve their primary score. We therefore
returned to embeddings as our main inputs. Our interpretation is that the
embeddings already describe the stopping reason directly, while the other
fields describe broader clinical context. With few observed labels, extra
sparse dimensions and missing drug annotations could add noise more easily
than useful information. These experiments supported a simpler representation,
although we did not repeat every feature screen on the final blend.

Random CV did not address our concern about predicting future trials: the test
trials started later, and the mix of stopping reasons could change over time.
We added earlier-to-later checks to see how models behaved under that shift.
They gave us another benchmark, but did not consistently identify the best
Kaggle candidate. We also investigated the combined evaluation score $J$,
defined in the benchmark notebook as $J=(A+B)/2$. Here $A$ is ten-repeat random
CV on 235 labels, and $B$ is the mean of three chronological checks. This
experiment used different split boundaries from our grouped benchmark and aimed
to rank models by broader generalization. Its rankings sometimes disagreed
with Kaggle. We learned that averaging two
plausible checks does not automatically produce a dependable selection metric;
we kept chronology as a separate diagnostic rather than relying on $J$ alone.

The per-reason results also showed that a strong overall score could hide
weaknesses in Funding and Other. We tried extra weighting for those labels
and nonlinear specialists for Other, hoping that a focused model could recover
cases the shared classifier missed. The weighting changes did not improve
decisions, and the specialists did not give a reliable overall gain. We also
found that almost all generated examples had only one positive reason,
suggesting that self-training could reinforce missed secondary labels.
Leaving unselected Other labels unknown recovered more true Other cases, but
also introduced enough false positives to reduce overall F1. The lesson was
that improving recall for a difficult reason was insufficient if precision
fell too far. Funding remained especially hard to assess with so few examples.

These findings shaped our blending strategy. We kept most of the weight on
the stronger anchored model and used LDA or Cand2 to contribute different
prediction patterns. Cand2's weaker standalone result did not rule out a
useful contribution, and its small share improved the final blend. However,
learned combiners, explicit diversity penalties, log-odds pooling, and symbolic
regression did not improve our selected result. Our strongest model combined different approaches without moving far from the
model that already worked well.

\subsection{Part 2: temporal shift and error trade-offs}

In Part~2, better log loss did not always mean better F1: a probability can
become more accurate while crossing the 0.5 threshold in a way that lowers F1.
Because the test trials start later and success rates rose within our labeled
data, we prioritized chronological evaluation. Recency
weighting and recent neighbors reduce the influence of older trials. The
earlier ensemble helped through different error patterns, although increasing
its share from 20\% to 30\% improved Kaggle log loss by only 0.00009. File~15's
worse leaderboard result reinforced the limits of historical averages.

Our error analysis suggests that Phase~III failures are particularly difficult.
These trials may reach completion but miss their endpoints, leaving eligibility
text similar to that of successful trials. Recruitment and funding failures
may reflect administrative constraints, and the ensemble identifies some
additional cases. This motivated adding clinical context beyond eligibility
text, but the final model still trades improvements in some groups for errors
in others. These
explanations remain tentative: Safety, Funding, and Other have few examples in
each phase, and predicted reason tags can place trials in the wrong category.
The no-reason groups contain trials absent from Part~1, rather than a predicted
``no reason'' category.

\smallskip
\noindent\textbf{AI disclosure.}
We used AI tools to assist with model improvements, resolve errors, automate
testing, and edit and format this report.

{
    \small
    \bibliographystyle{ieeenat_fullname}
    \bibliography{main}
}

\end{document}
